% Figaro Paper F: The Wallet as Legal Subject
% Subordination, Property, and the Dissolution of the Employee–Contractor Line
% Audience: labor-law scholars, legal philosophers, COALA-adjacent
% practitioners, gig-economy policy researchers.

\documentclass[11pt,a4paper]{article}

\usepackage[utf8]{inputenc}
\usepackage[T1]{fontenc}
\usepackage{amsmath,amssymb,amsthm}
\usepackage{mathtools}
\usepackage{booktabs}
\usepackage{graphicx}
\usepackage{hyperref}
\usepackage{cleveref}
\usepackage{enumitem}
\usepackage{xcolor}
\usepackage[margin=1in]{geometry}
\usepackage{natbib}
\bibliographystyle{plainnat}

\theoremstyle{remark}
\newtheorem{remark}{Remark}[section]

\newcommand{\figaro}{\textsc{Figaro}}

\title{%
  The Wallet as Legal Subject:\\
  Subordination, Property, and the Dissolution\\
  of the Employee--Contractor Line%
}

\author{
  Alessandro Daliana\thanks{Corresponding author. \texttt{adaliana@figaro.org}}
}

\date{April 2026}

\begin{document}

\maketitle

% ════════════════════════════════════════════════════════════════════
\begin{abstract}
Labor law distinguishes employees from independent contractors along
a single doctrinal axis: the degree of subordination the employer
exercises over how the worker conducts themselves within the
employer's system of processes and values. The axis is a continuum,
the classification is a threshold placed somewhere on it, and the
threshold is the perpetual subject of legislative and judicial
contest --- AB5 and Proposition 22 in California, \emph{Uber BV v
Aslam} in the United Kingdom, the 2024 EU Platform Work Directive.
The cases are all the same case, relitigated in each jurisdiction
against each new technology.

This paper argues that the bonded settlement primitive
of the mechanism paper, implemented and verified
in the implementation paper, changes the doctrinal terrain by
eliminating the axis. Figaro parties are bonded counterparties;
there is no employment relation to classify. The wallet collapses a
distinction the Roman legal tradition kept separate: \emph{res}
(property) and \emph{persona} (legal party). Every productive asset
is simultaneously property and party --- self-owned rather than
owned or rented.

We develop the argument in three moves. First, we locate
subordination as the doctrinal variable in the long-running labor
classification debate, drawing on Anderson's analysis of the firm
as private government~\citep{anderson2017private}, Ellerman's
argument that the employment contract renews the master--servant
relation in softer form~\citep{ellerman1992property}, and Taleb's
compressed version of the same point in \emph{Skin in the
Game}~\citep{taleb2018skin}. Second, we show that the gig-economy
case law is a classification dispute whose ground disappears once
subordination ceases to be a variable. Third, we name the
trade-off: Figaro is a leveller on the subordination axis but
simultaneously a privatizer of employment-attached protections
(unemployment insurance, healthcare, pensions, collective
bargaining). What follows is a matter of policy, not mechanism.
\end{abstract}

\medskip
\noindent\textbf{Keywords:}
labor law, employee classification, gig economy, private government,
subordination, legal subject, self-sovereign identity

% ════════════════════════════════════════════════════════════════════
\section{Introduction}
\label{sec:intro}

Labor law, across every jurisdiction that has one, answers a
recurring question: is this person an employee or an independent
contractor? The answer matters because the two categories attract
radically different legal consequences --- minimum wage, overtime,
unemployment insurance, employer health-insurance obligations,
workers' compensation, collective-bargaining rights, anti-discrimination
protections, tax treatment, pension eligibility. The classification
test varies by jurisdiction, but all tests converge on a common
underlying variable: the degree to which the putative employer
directs and controls the putative employee's work --- their
schedule, their tools, their processes, their supervision, their
ability to work for others.

We will call this variable \emph{subordination}, following the
usage common in Continental labor-law scholarship and in the
political-economy tradition we return to in \Cref{sec:subordination}.
Subordination is a continuum; the employee--contractor line is a
threshold placed somewhere on it; and the location of the threshold
is the perpetual subject of doctrinal and political contest.

The gig economy made the contest legible. Drivers, couriers, and
platform freelancers occupy a region of the subordination continuum
the existing legal categories were not designed for: loosely
supervised by algorithm, nominally free to refuse work, economically
dependent on the platform nonetheless. Jurisdictions have responded
with increasingly granular tests. California's Assembly Bill 5
(2019) codified the ABC test~\citep{ab5_2019}, pushing many gig
workers toward employee status; Proposition 22 (2020) carved out an
exception for app-based rideshare and delivery
drivers~\citep{prop22_2020}. The UK Supreme Court in \emph{Uber BV
v Aslam}~\citep{uber_aslam_2021} classified drivers as ``workers''
--- a third category intermediate between employee and contractor
--- on subordination grounds. The EU Platform Work Directive
(Directive (EU) 2024/2831)~\citep{eu_platform_directive_2024}
established a presumption of employment for platform workers
meeting specified control indicators.

Each of these cases and statutes is, at root, an attempt to draw
the same line differently. The line itself is what the cases
cannot dispense with.

\paragraph{The argument of this paper.}
We argue that the bonded settlement primitive developed
in the mechanism paper and implemented
in the implementation paper changes the question. The companion
institutional-economics paper develops the
institutional consequence: the firm becomes one organizational
pattern among several when the subordination overlay is no longer
enforcement-driven. The legal consequence follows but needs a
separate treatment because its vocabulary and its audience are
different.

Specifically: when every party to a productive exchange is a bonded
counterparty --- signing commitments under their own key, posting
collateral against their own performance, receiving settlement
directly to their own wallet --- there is no employment relation
to classify. The classification problem has no subject. The
employee--contractor line is not redrawn; it is dissolved by the
absence of the relation it was drawing.

\paragraph{What this paper is not.}
This paper does not provide legal advice. It does not argue that
labor law should be abolished. It does not argue that
employment-attached protections are undesirable. It does not
predict that the bonded primitive will displace employment as the
dominant coordination form, and does not endorse any particular
outcome of the classification it dissolves. Its argument is
diagnostic: the primitive changes the doctrinal terrain, and
legal scholarship should understand that change independently of
any normative conclusion one might draw from it.

\paragraph{Paper organization.}
\Cref{sec:subordination} develops subordination as the
labor-law variable, drawing on Anderson, Ellerman, and Taleb.
\Cref{sec:primitive} summarizes the settlement primitive in the
terms relevant here. \Cref{sec:wallet} develops the central legal
move: the wallet as a legal subject that collapses the classical
\emph{res}/\emph{persona} distinction. \Cref{sec:gig} re-reads the
gig-economy case law through that lens.
\Cref{sec:leveller} names the leveller/privatizer trade-off.
\Cref{sec:conclusion} concludes.

% ════════════════════════════════════════════════════════════════════
\section{Subordination as the Labor-Law Variable}
\label{sec:subordination}

\subsection{The Doctrinal Test}

Every mature labor-law regime operationalizes the
employee--contractor distinction through a control-flavored test.
The details vary:

\begin{itemize}
  \item The IRS twenty-factor test in the United States examines
    behavioral control, financial control, and the nature of the
    relationship.
  \item California's ABC test, codified in AB5, places the burden
    on the hiring entity to show that the worker
    (A) is free from control and direction,
    (B) performs work outside the hiring entity's usual business,
    and (C) is customarily engaged in an independently established
    trade.
  \item UK jurisprudence, following \emph{Autoclenz}~\citep{autoclenz_2011}
    and refined in \emph{Aslam}, examines the reality of the
    working arrangement rather than its contractual label,
    focusing on personal service, mutuality of obligation, and
    control.
  \item Continental systems (France, Germany, Italy) distinguish
    \emph{travail salarié}, \emph{abh\"angige Arbeit},
    \emph{lavoro subordinato} respectively, anchoring the status
    on \emph{subordination juridique} --- the juridical subordination
    of the worker to the employer's directive power.
\end{itemize}

The tests differ in which facts they weight and in whether the
burden sits on the worker or the employer. They agree on what the
variable is. Control, direction, supervision, integration into the
employer's processes --- these are all labels for the same underlying
relationship: subordination of the worker's conduct to the employer's
authority, in exchange for a wage.

\subsection{The Philosophical Position}

A scholarly tradition with roots in classical liberalism and in
Marx, refreshed more recently by legal and political philosophers,
reads the employment relationship as lying on the same continuum
as older forms of property-in-persons. We summarize three
representative positions; the paper does not endorse any of them,
and each can be read as descriptive rather than normative.

\paragraph{Ellerman: the firm as rent-a-person.}
\citet{ellerman1992property} argues that the employment contract is
not merely economically one-sided but ontologically incoherent: it
attempts to rent \emph{persona} (the worker as legal party) on the
same terms as \emph{res} (the worker's labor power as object). In
ordinary jurisprudence, persons cannot be rented --- a doctrine
that predates and survived the legal abolition of chattel slavery.
The employment contract, Ellerman argues, honors the formal
doctrine by contracting for labor power rather than the person, but
the substance of the relation (subordination of the worker's
conduct to another's authority for a defined period) exercises the
same operative control. The argument is philosophical rather than
legal --- Ellerman is not claiming that employment contracts are
unenforceable --- but it names a tension that labor law
operationalizes as the employee/contractor test.

\paragraph{Anderson: the firm as private government.}
\citet{anderson2017private} makes the complementary argument at the
level of political theory. Firms exercise governmental powers over
employees --- mandatory schedules, speech restrictions, body
surveillance, drug testing, arbitrary dismissal --- that are
constitutionally proscribed when exercised by the state. American
political vocabulary, Anderson argues, has no term for this because
the liberal tradition's attention to state power left the
authoritarian structure of the workplace unanalyzed. The firm is
a private government, and the employment relation is the
governance relation within it.

The structural observation common to Ellerman and Anderson --- that
the employment relation places one party's conduct under another
party's authority along an axis whose endpoints touch
historical arrangements of chattel and serfdom --- has been put
in popular form by other writers as well.\footnote{See, e.g.,
\citet[][Book 4, Ch.\ 3, ``How to Legally Own Another
Person'']{taleb2018skin}, who frames the relation as ``an
employee is practically a slave.'' Taleb's register is rhetorical
and his treatment is not peer-reviewed; we cite him here as
evidence that the observation has currency outside academic labor
philosophy, not as an independent argument on par with Anderson
and Ellerman. The substantive analytical work in this paper rests
on the latter two.} The observation, however framed, is that the
axis is subordination; the differences between historical
arrangements at different points on it are quantitative. The
paper does not endorse this framing as a moral claim about
modern employment; we return to that scope question in §7.

\paragraph{Subordination in two registers.}
The doctrinal test in \Cref{sec:subordination} and the
philosophical treatment just given describe the same
relationship from two disciplinary vantages. The doctrinal register
is direction-and-control: does the putative employer specify
schedule, tools, processes, supervision, ability to work for
others? The property-rights register, developed in the
companion institutional-economics paper (§3.1),
is residual-rights allocation: does the employment contract
transfer residual control rights over non-contractible decisions
to an ownership nucleus separate from the productive asset?
These are not competing definitions. The direction-and-control
test is the \emph{operational} face of subordination --- what a
court observes when assessing the relation. The residual-rights
framing is its \emph{structural} face --- what a property-rights
theorist sees when analyzing the same relation. The two
converge: a relation with high directional control is a relation
in which residual rights have been transferred; a relation in
which residual rights stay with the worker is a relation of low
directional control. This paper uses the doctrinal register
(because the subject is labor law); the companion paper uses the
property-rights register (because the subject is the boundary of
the firm). The primitive dissolves both faces at once because
they are faces of the same thing.

\subsection{Why the Threshold Keeps Moving}

Once subordination is recognized as the variable and the
employee--contractor line as a threshold placed on it, the
perpetual legal contest becomes explicable. The threshold has no
principled location. Any specific test produces a bright line
through a continuum, and whichever side of the line a given
working arrangement falls on becomes the object of rent-seeking
effort by the party for whom that side is cheaper.

\begin{itemize}
  \item Employers push the threshold upward: reclassifying employees
    as contractors reduces payroll taxes, avoids
    benefits-attachment, eliminates collective-bargaining
    obligations, and removes discrimination-law exposure. Every
    major gig-economy platform has been the subject of litigation
    on this point.
  \item Workers with bargaining power push the threshold downward:
    reclassification as employees unlocks minimum-wage
    protections, overtime, employer-funded health insurance,
    unemployment eligibility, and union rights. Much of AB5 was a
    legislative response to this pressure.
\end{itemize}

The threshold is not ``wrong''; it is \emph{under-specified}. The
underlying continuum has no natural joint. As technologies of
remote direction (apps, algorithms, performance dashboards)
evolve, the reality of the control relation changes faster than
the law can locate the threshold, which is why gig-economy cases
recur across jurisdictions with broadly similar facts and broadly
divergent outcomes.

% ════════════════════════════════════════════════════════════════════
\section{The Settlement Primitive}
\label{sec:primitive}

We summarize only the features relevant to the legal argument. The
mechanism is established in the mechanism paper; the
Solidity implementation and formal verification
in the implementation paper; the evidentiary framework for off-chain
dispute adjudication in the evidence-and-law paper.

\paragraph{Bonded counterparties.}
A transaction between two parties is consummated by dual-signed
EIP-712 commitments. Each party signs from their own wallet. Each
party locks collateral against their own performance --- the buyer
locks $2P$, the seller locks $2G$ where $G$ is the cumulative
upstream value bonded in the process. Only the buyer can trigger
resolution. Defection is structurally more expensive than
cooperation. There is no platform between the parties, no
arbitrator authorized to resolve disputes on-chain, and no admin
key permitting external override.

\paragraph{What the primitive does not require.}
No employment relation, no subcontractor hierarchy, no master--servant
agreement, no W-9 or W-2, no 1099 or P-60, no PAYE code, no
\emph{contrat de travail}. The primitive operates on parties as
parties: each is a wallet address bonded against performance for a
specified exchange. The direction-and-control relation that labor
law categories turn on does not exist at the primitive layer; it is
replaced by the direction given by the off-chain agreement both
parties sign and the bond that backs their mutual performance.

% ════════════════════════════════════════════════════════════════════
\section{The Wallet as Legal Subject}
\label{sec:wallet}

Roman private law distinguished \emph{res} (things, property,
objects that can be owned and alienated) from \emph{persona}
(persons, parties, subjects capable of acting legally). The
distinction survived and structures modern private law: physical
assets are \emph{res}; natural persons and recognized legal
persons (corporations, states, partnerships) are \emph{personae};
the two interact through contracts of sale, lease, license, and
labor.

The employment contract sits awkwardly in this structure. A worker
is unambiguously a \emph{persona}, but the employment contract
treats the worker's labor power as \emph{res} --- a resource
rented from the worker for a specified period, under the
employer's direction. Ellerman's argument
in~\Cref{sec:subordination} is that this treatment is operationally
coherent only because the legal tradition has learned to overlook
the \emph{persona}/\emph{res} boundary in this one setting. Every
labor-law test for employee status is implicitly an attempt to
police that boundary without naming it.

\paragraph{The wallet as collapse.}
A wallet in the \figaro{} architecture is a legal artifact that
does not fit the classical categories. It is a property --- a
cryptographic key-pair that its holder owns and can transfer or
abandon. It is simultaneously a party --- it signs commitments,
posts bonds, receives settlements, accrues credentials, and
generates an evidence record that off-chain adjudication can
consume (evidence-and-law paper). The wallet is the productive
asset acting in its own legal name.

The collapse is not metaphorical. When the party to a commitment
is the wallet itself (whether operated by a natural person, a
software agent, or an organizational treasury), the classical
\emph{persona}/\emph{res} partition no longer partitions. The
wallet is what it controls and it controls what it is.

\paragraph{Three consequences.}
\begin{enumerate}[label=(\alph*)]
  \item \textbf{Self-control is the default.} A wallet cannot be
    \emph{owned} by another wallet in the way a \emph{res} can be
    owned by a \emph{persona}: ownership transfer requires key
    transfer, which is a change of control rather than a change of
    subject. We are careful here not to collapse ownership and
    control: the property-rights tradition (Grossman--Hart, Hart--Moore)
    builds on the distinction between the two, and the present
    paper relocates the question rather than dissolving it. The
    relocation is that residual control rights over the productive
    capability the wallet represents are held by whoever holds the
    key, not by an ownership nucleus separate from the productive
    asset. \emph{What} the key controls --- the cryptographic act
    of signing, but also indirectly the deployment of physical
    assets, tacit knowledge, and judgment under uncertainty that
    the signer brings to bear --- remains a property-rights
    question with the same structure as before; it is just answered
    differently.
  \item \textbf{No rental of persona.} A commitment between two
    wallets bonds each wallet's performance directly. There is no
    construct in which one wallet's persona is rented to another
    for a defined period; each party signs for themselves,
    commitment-by-commitment.
  \item \textbf{Credentialing is self-attached.} At the protocol
    tier (per the kernel/protocol/runtime distinction
    in §1.4 of the mechanism paper), schema-validated
    attestations and operator profiles attach to the wallet; at
    the runtime tier, NFT-anchored deeds and credentials may be
    layered on. A worker's skills, certifications, and
    performance history travel with the wallet, not with the
    employer. The resume becomes the wallet; the reference becomes
    an attestation.
\end{enumerate}

These consequences are descriptive. They do not depend on any
policy choice; they follow from the architecture of dual-signed
bonded commitments settled directly between parties.

\paragraph{A fourth consequence: the wallet is self-accounting.}
The Roman private-law inheritance distinguishes \emph{res}
from \emph{persona} and keeps both separate from the
\emph{liber rationum} --- the book of account in which the
entity's dealings are recorded. A \figaro{} wallet collapses
the first two (§4 above); it also collapses both into the
third. The companion accounting
paper develops the argument
that a \figaro{} process is a self-closing ledger period
whose records are constituted by the protocol rather than
maintained by any party. For labor-law analysis, the
implication is that the wallet is not only property-and-party
but also its own continuously updated book of account: the
worker-as-wallet has a record of performance, attested
capability, and settled compensation that travels with the
wallet indefinitely. This eliminates one traditional source
of asymmetry between employer and employee --- the employer
keeps the books, the employee does not --- by making the
books protocol-level and symmetrically readable.

% ════════════════════════════════════════════════════════════════════
\section{Displaced and Stateless Subjects}
\label{sec:displaced}

The arguments of \Cref{sec:wallet} presumed a conventional
institutional background: persons with stable civil-legal
subjecthood within a jurisdiction, interacting through (or
outside) the employment contract. The wallet-as-legal-subject
argument compresses two classical categories for readers
embedded in that background. For a second class of readers ---
populations whose access to stable civil-legal subjecthood is
itself disrupted --- the same argument discloses a different
consequence worth stating separately.

\paragraph{Scale.}
The United Nations Refugee Agency estimated 120 million forcibly
displaced persons worldwide at the end of 2023, and approximately
4.4 million persons registered as stateless, with a substantially
larger estimated population whose statelessness is undocumented
\citep{unhcr2024globaltrends}. An overlapping population lacks
stable banking and enforcement access for reasons of informal
migration status, residence without papers, or residency in
jurisdictions whose financial infrastructure does not reach
them. The population for whom the labor-law classification debate
in \Cref{sec:subordination} is a second-order question --- because
the first-order question is whether they are recognized parties to
any contract in any forum --- is substantial.

\paragraph{The received architecture.}
Labor law, contract law, and commercial-dispute adjudication each
presume a recognized legal subject attached to a jurisdiction that
enforces their rights. Refugees, stateless persons, and
undocumented migrants occupy the negative space of that
presumption. The employment relation is typically unavailable to
them formally (permission-to-work requirements), the
independent-contractor relation unstable (cannot register a
business, cannot enforce contracts in the local courts), and
banking access restricted or absent. The humanitarian-economics
literature on refugee economies~\citep{betts_collier_2017} and the
policy frontier between camp and labor
market~\citep{pritchett2006people} documents the economic cost of
this gap: the productive capacity of the displaced population is
systematically underused because the institutional scaffold for
engaging it is absent.

\paragraph{What the primitive does and does not do.}
The primitive does not resolve statelessness. It does not confer
citizenship, restore legal recognition, or create enforcement
against a state that denies the person's standing. It is not a
substitute for the political work that addresses those gaps.
What it does is narrower: it provides a commerce architecture
whose precondition is a cryptographic key rather than a recognized
civil-legal subjecthood. A wallet does not require a passport to
bond; a bonded commitment does not require a host-state court to
enforce; a settled exchange does not require a domestic bank to
clear. The primitive opens a region of productive activity ---
bilateral, bonded, settled between strangers --- to participants
for whom the legal scaffold that conventionally underwrites such
activity is unavailable.

\paragraph{Arendt and the right to have rights.}
\citet{arendt1951origins} observed that the fundamental human
right is the right to have rights --- the right to belong to a
political community within which one's other rights have
standing. Her analysis of statelessness named the condition as a
civilizational failure rather than an individual misfortune. The
bonded primitive does not grant the right to have rights; it does,
more narrowly, grant a capacity to have commerce. The two should
not be confused. The right to have rights is political and is not
within the gift of any technology; the capacity to have commerce
is structural and becomes operational once the cryptographic
infrastructure is in reach. For populations for whom the first
right is denied, the second capacity is non-trivial.

\paragraph{A cautious posture.}
A paper in labor-law scholarship is not the venue to develop the
humanitarian consequences of the primitive in detail; the brief
treatment here is positional rather than programmatic. We note two
constraints. First, key custody is a real constraint: populations
with unreliable access to power, devices, and connectivity face
usability barriers at the cryptographic layer that are not solved
by the primitive per se. Second, counterparty willingness: a
well-resourced buyer may decline to bond with a counterparty
lacking civil-legal standing for reasons that have nothing to do
with the primitive and everything to do with downstream legal
exposure in the buyer's home jurisdiction. These constraints are
empirical and evolve with the infrastructure. The analytical
point is that the primitive moves the humanitarian-economics gap
from an architectural impossibility to an operational engineering
question, which is a nontrivial move.

% ════════════════════════════════════════════════════════════════════
\section{The Gig-Economy Cases, Re-read}
\label{sec:gig}

The gig-economy case law can be read as a series of attempts to
place the employee--contractor threshold in the face of a
technology (the platform) that had engineered a new position on the
subordination continuum. We re-read four landmark cases/statutes
through the primitive's lens.

\paragraph{California AB5 (2019).}
AB5 codified the ABC test established in \emph{Dynamex Operations
West v Superior Court}~\citep{dynamex_2018}, shifting the
classification default toward employee status for workers
performing services in the hiring entity's usual course of
business. The statute's legal force turns on the direction-and-control
finding (prong A), the nature-of-work finding (prong B), and the
independent-trade finding (prong C).

Under the primitive, none of the three prongs are well-posed. There
is no hiring entity to direct and control (prong A has no subject).
There is no usual course of business for the bonded counterparty
relation (prong B presupposes a firm). The bonded party is by
construction independently established (prong C is trivially
satisfied, not by design but by default). AB5 operationalizes a
threshold that Figaro does not instantiate at all.

\paragraph{Proposition 22 (2020).}
Prop 22 created a statutory exception to AB5 for app-based
rideshare and delivery drivers, classifying them as contractors
with a specified minimum earnings guarantee and healthcare subsidy.
The political coalition behind Prop 22 argued that AB5 was forcing
classification into a category inappropriate for the work; the
coalition against argued that the exception hollowed out AB5's
protective purpose. Both arguments are downstream of the threshold
choice.

Under the primitive, Prop 22's carve-out is unnecessary because
the category it is carving out of does not apply. The underlying
question it tries to answer --- how to provide minimum earnings
and healthcare to workers outside the employment relation --- is
what we treat in \Cref{sec:leveller} as the leveller/privatizer
trade-off. Figaro does not answer it; it changes what the question
is.

\paragraph{Uber BV v Aslam (UKSC 2021).}
The UK Supreme Court held unanimously that Uber drivers are
\emph{workers} --- an intermediate category in UK employment law
that attracts minimum wage and holiday-pay protections but not
full employee rights. The Court examined the substance of the
relation (algorithmic direction, tightly specified routes, rating
systems, ability-to-work-for-others constraints) rather than the
contractual label. The judgment is unusually clear that
subordination-in-substance is the operative test.

Under the primitive, the Aslam analysis does not run. There is no
entity whose substance can be examined against the contractual
label, because the primitive produces no contract between
substance-behind-label and worker. A bonded counterparty signs
commitment-by-commitment with the party it serves. The
rating-and-algorithmic-direction apparatus that Aslam turned on
is a platform artifact; it does not exist in the primitive.

\paragraph{EU Platform Work Directive (2024).}
The Directive establishes a presumption of employment for platform
workers when specified criteria are met (algorithmic management,
remuneration caps, appearance restrictions, tool prescriptions,
supervision of work, restrictions on working for others). Member
states are to implement by 2026. The Directive is the most
sophisticated of the four measures: it treats platform employment
as its own typological object, with a presumption rule rather than
a bright-line test.

Under the primitive, the Directive's presumption criteria have no
referent --- there is no platform between parties whose algorithmic
management would trigger the presumption. The Directive applies to
the platform-mediated region of the contemporary labor market;
Figaro moves exchanges out of that region entirely.

\paragraph{The pattern.}
Each measure engages the subordination variable directly, and each
would be either vacuous or inapplicable if the parties to an
exchange were bonded counterparties rather than parties to a
platform-mediated or firm-mediated employment relation. The cases
do not fail in the Figaro setting; they have no application there.

% ════════════════════════════════════════════════════════════════════
\section{Leveller and Privatizer}
\label{sec:leveller}

The primitive is, in one register, a leveller. It eliminates a
structural asymmetry in labor markets: the subordination gradient
between employer and employee. Every party is bonded on the same
terms; every party signs from its own wallet; no party is rented
as \emph{persona}.

The primitive is, in another register, a privatizer. The protections
labor law built around the employment relation --- employer-funded
healthcare, employer-paid payroll taxes feeding unemployment
insurance, employer pension contributions, workplace safety
inspections tied to employer status, collective bargaining
structured around bargaining units within firms --- are all
attached to the category that the primitive removes. When
employment is not the container, employment-attached protections
are not delivered.

\paragraph{Asset specificity is a scope condition.}
The argument above presumes that the productive activity in
question is the kind for which symmetric bilateral bonding is an
adequate substitute for hierarchical governance. The
companion institutional-economics paper (§3) is
explicit that asset specificity continues to favor hierarchy:
when a productive asset's value is contingent on a specific
counterparty (relationship-specific human capital, co-located
equipment, firm-specific tacit routine), the holdup risk lies in
the pre-transaction sunk investment rather than in the
transaction itself, and symmetric bonding does not neutralize
that holdup --- it relocates it. The class of activity for which
the primitive substitutes for the firm is therefore bounded.
What class is bounded by what level of specificity is an
empirical question outside this paper's scope. The reader should
take F's argument as applying within the bound: where asset
specificity is low or where specificity-protecting arrangements
are achievable through other means (long-term peer networks,
treasury communities, serial bonded-commitment patterns), the
employment relation becomes optional. Where it is high, the
employment relation persists for the standard TCE reason.

\paragraph{Four options for the political response.}
We do not advocate for any of them; we name them so the trade-off
is legible.

\begin{enumerate}[label=(\roman*)]
  \item \textbf{Attach protections to the wallet.} Health
    insurance, pension contributions, income protection become
    benefits that the wallet itself purchases from portable
    providers, priced per unit of bonded activity. This is the
    direction Prop 22 gestured at, without the AB5 backdrop
    making it necessary.
  \item \textbf{Attach protections to citizenship.} Universal
    basic income, single-payer healthcare, state-funded
    unemployment assistance --- all decouple protection from
    employment and re-couple it to polity membership. The social
    democratic tradition has advocated this independently of
    Figaro; the primitive removes one historical argument against
    the decoupling (the argument that employment was the natural
    locus because that is where most people spend most of their
    time).
  \item \textbf{Preserve the employment option.} Nothing in the
    primitive prevents the employment relation from continuing to
    exist. Firms can hire employees, employees can bargain
    collectively, labor law can continue to regulate the
    subordination region. The primitive coexists with the
    employment relation rather than abolishing it. The question
    is what fraction of productive activity moves out of that
    region, and at what rate.
  \item \textbf{Treat employment as the asset-specificity
    residual.} A position adjacent to (iii) but more pointed:
    employment continues to apply to the class of productive
    activity for which Williamsonian asset specificity makes
    hierarchical coordination the equilibrium response, and
    labor-law protections continue to attach to that residual
    domain. On this view the primitive does not compete with
    employment in the high-specificity region at all; it
    competes only in the low-specificity region, where
    employment was a second-best response to enforcement cost
    rather than a first-best response to specificity. The
    question becomes empirical: what is the asset-specificity
    distribution of contemporary productive activity, and what
    fraction lies on each side of the threshold? F does not
    answer the question; we name it because the institutional-economics
    tradition would otherwise treat its absence as oversight.
\end{enumerate}

The trade-off is real: each of (i)--(iv) has different
distributive consequences, and the fight over which one
dominates will occupy the political economy of the next several
decades regardless of how the primitive diffuses. What the
primitive changes is that the fight becomes explicit: the
attachment of protections to employment was a historical
contingency that looked like a natural fact; once the category
is removable, the attachment becomes a policy choice.

% ════════════════════════════════════════════════════════════════════
\section{Conclusion}
\label{sec:conclusion}

Labor law has treated subordination as an axis on which to place
threshold after threshold, in each jurisdiction against each
technology. The bonded settlement primitive removes the axis by
removing the relation that gave it its subject matter. The
employee--contractor line, the platform-worker line, the \emph{worker}
category, the \emph{subordonné} test --- all presuppose a
directional control relation between two parties that the primitive
does not produce. The classification problem has no classification
problem.

What is left is a different set of questions. What protections, if
any, should attach to productive activity independent of the
employment relation? Which organizational patterns will absorb the
productive activity that moves out of the employment relation, and
how will they be regulated? How will labor law's traditional
counterparty, the state, relate to productive activity that
increasingly occurs between wallets rather than between firms and
workers? These questions are not answered here and are not, in our
reading, answerable from mechanism-design first principles. They
are legal-doctrinal, political-economic, and constitutional
questions that the primitive poses rather than resolves.

\paragraph{A scope caveat.}
This paper's argument is diagnostic. The primitive changes what
labor law is asked about; the paper says so. The paper does not
argue that the primitive will diffuse uniformly or that employment
as a legal category will disappear. It does not argue that the
disappearance of the classification problem is desirable, in whole
or in part. It argues that the doctrinal terrain has changed,
that legal scholarship has not yet caught up to the change, and
that whatever normative position the reader takes on the change
will be better served by engagement with the primitive on its own
terms than by treating it as one more gig-economy platform to be
reclassified under the existing test.

\paragraph{A note on rhetoric.}
We have, following Anderson and Ellerman (and, in a popular
register, Taleb), used the vocabulary of ``subordination'' and
``slavery'' to describe the employment relation. The rhetorical
register is deliberate but should not be mistaken for a normative
claim about modern employment. Modern labor law --- at-will
doctrine carve-outs notwithstanding --- has developed substantial
protections that distinguish the employee from the serf and the
serf from the slave. The axis on which these arrangements sit is
structural, not moral; the primitive removes the arrangement, not
the moral progress the arrangement's softening represented.
Whether that removal is desirable is a policy question; whether
it is happening is an engineering one; and the two questions
should not be conflated.

\paragraph{A cross-cluster note on the diagnostic posture.}
This paper, along with the companion institutional-economics
paper, the companion political-economy paper, and the companion
evidence-and-law paper, each individually disclaims
the normative register. Each is explicit that it is diagnostic
rather than predictive or prescriptive. Read together, the four
diagnoses nonetheless accumulate into a picture: a world in which
the firm demotes, gig-economy classification dissolves, platform
hegemony becomes visible, and legal adjudication cleanly decouples
from settlement. A careful reader can observe that four carefully
bounded diagnoses, read in sequence, approach an \emph{implicit}
endorsement --- an endorsement by exhaustion --- that no
individual paper makes. We wish to name this effect without
resolving it. The portfolio describes what it observes about the
primitive's consequences; it does not predict which of those
consequences materialize, does not advocate for those that do, and
does not compose the four diagnoses into a prediction. The
normative question of whether any of the diagnosed changes are
desirable is policy and politics, and the present paper --- along
with its companions --- leaves that question open.

\section*{Acknowledgments}
I am grateful to readers who pushed me to separate the legal
consequence of the settlement primitive from the institutional and
political-economy consequences developed in companion papers. The
Coalition of Automated Legal Applications (COALA) has been a
consistent interlocutor on the legal-doctrinal questions treated
here, and the Stanford CodeX community on the computational-law
side. The provocation of the title comes from Nassim Taleb's
chapter heading; the philosophical spine comes from Elizabeth
Anderson and David Ellerman, to both of whom my debt is evident
throughout.


% ════════════════════════════════════════════════════════════════════
% References
% ════════════════════════════════════════════════════════════════════

\begin{thebibliography}{99}

\bibitem[AB5(2019)]{ab5_2019}
California State Legislature.
\newblock Assembly Bill No. 5, Ch. 296, Stats. 2019.
\newblock Codifying the ABC test for employee classification.

\bibitem[Anderson(2017)]{anderson2017private}
Anderson, E.
\newblock \emph{Private Government: How Employers Rule Our Lives
  (and Why We Don't Talk About It)}.
\newblock Princeton University Press, 2017.

\bibitem[Arendt(1951)]{arendt1951origins}
Arendt, H.
\newblock \emph{The Origins of Totalitarianism}.
\newblock Harcourt, Brace \& Co., New York, 1951.
\newblock See especially Part II, Ch.~9 (``The Decline of the
  Nation-State and the End of the Rights of Man''), for the
  formulation of the right to have rights.

\bibitem[Betts and Collier(2017)]{betts_collier_2017}
Betts, A. and Collier, P.
\newblock \emph{Refuge: Transforming a Broken Refugee System}.
\newblock Penguin, London, 2017.

\bibitem[Autoclenz(2011)]{autoclenz_2011}
\emph{Autoclenz Ltd v Belcher}, [2011] UKSC 41.

\bibitem[Dynamex(2018)]{dynamex_2018}
\emph{Dynamex Operations West, Inc. v Superior Court},
  4 Cal.5th 903 (2018).

\bibitem[Ellerman(1992)]{ellerman1992property}
Ellerman, D.~P.
\newblock \emph{Property and Contract in Economics: The Case for
  Economic Democracy}.
\newblock Blackwell, Cambridge, MA, 1992.

\bibitem[EU Platform Work Directive(2024)]{eu_platform_directive_2024}
European Parliament and Council.
\newblock Directive (EU) 2024/2831 of 23 October 2024 on improving
  working conditions in platform work.
\newblock Official Journal of the European Union, 2024.

\bibitem[Pritchett(2006)]{pritchett2006people}
Pritchett, L.
\newblock \emph{Let Their People Come: Breaking the Gridlock on
  Global Labor Mobility}.
\newblock Center for Global Development, Washington, DC, 2006.

\bibitem[Prop 22(2020)]{prop22_2020}
California Proposition 22 (2020), codified at Cal. Bus. \&
  Prof. Code §§ 7448--7467.

\bibitem[Taleb(2018)]{taleb2018skin}
Taleb, N.~N.
\newblock \emph{Skin in the Game: Hidden Asymmetries in Daily Life}.
\newblock Random House, New York, 2018.
\newblock See especially Book~4 (``Wolves Among Dogs''), Chapter~3
  (``How to Legally Own Another Person'').

\bibitem[Uber v Aslam(2021)]{uber_aslam_2021}
\emph{Uber BV and others v Aslam and others}, [2021] UKSC 5.

\bibitem[UNHCR(2024)]{unhcr2024globaltrends}
United Nations High Commissioner for Refugees.
\newblock \emph{Global Trends: Forced Displacement in 2023}.
\newblock UNHCR, Geneva, 2024.

\end{thebibliography}

\end{document}
